\section{Moduł z licznikiem Geigera-Müllera}
\label{section:geiger_code}
Obsługa licznika odbywa się za pośrednictwem struktury \texttt{GeigerCounter}. Sprzętowy licznik T1 mikrokontrolera skonfigurowany jest do zliczania impulsów zewnętrznych generowanych przez lampę [AI mówi rurę] Geigera-Müllera. Co sekundę zliczona wartość jest odczytywana i zapisywana do 60-elementowego bufora cyklicznego.

Funkcja \texttt{cpm} zwraca liczbę zliczeń na minutę poprzez ekstrapolację zsumowanych wartości bufora na pełną minutę, nawet w przypadku niepełnego zapełnienia bufora. Takie podejście ma dwie konsekwencje: przy czasie pracy krótszym niż minuta wynik może być obarczony większym błędem, natomiast przy pełnym buforze nagłe skoki wartości są wygładzane przez uśrednianie.

Kod odpowiedzialny za odczyt i reset licznika T1 przedstawiono na listingu~\ref{lst:read_geiger_counter}. W celu zabezpieczenia przed przerwaniem wykonywania kodu zastosowano funkcję \texttt{interrupt::free} z biblioteki \texttt{avr\_device}.

\begin{lstlisting}[language=Rust, style=colouredRust, caption={Odczyt zliczeń licznika T1 mikrokontrolera},
                   label={lst:read_geiger_counter}]
fn timer1_read_and_reset(tc1: &TC1) -> u16{
    avr_device::interrupt::free(|_| unsafe {
        let val = tc1.tcnt1().read().bits();
        tc1.tcnt1().write(|w| w.bits(0));
        val
    })
}
\end{lstlisting}